\section{Kesesatan Berpikir Formal (Logical Fallacies)}

Selain argumen valid, terdapat pola \textbf{inferensi invalid} yang sering mengecoh. Dua fallacy formal yang menyerupai Modus Ponens dan Modus Tollens—namun tidak sah—disebut \textit{Fallacy} \cite{rosen2019,iep_natded,libretexts_proofs}.

\subsection{1. Menetapkan Konsekuen (Fallacy of Affirming the Consequent)}

Fallacy ini terjadi ketika seseorang keliru menganggap implikasi ($p \rightarrow q$) bersifat dua arah (bikondisional). 

\begin{teorema}[Jebakan Menetapkan Akibat]
Dari implikasi ($p \rightarrow q$) dan fakta bahwa konsekuen ($q$) benar, seseorang keliru menyimpulkan anteseden ($p$) benar. \textbf{Argumen tidak sah (invalid)}. 
\end{teorema}

\textbf{Anatomi Argumen Sesat:}
\[
\begin{array}{cl}
p \rightarrow q & \text{(Premis Mayor)} \\
q & \text{(Sang Akibat justru digandakan/dibenarkan)} \\
\hline
\therefore p & \text{\xcancel{(KESIMPULAN CACAT! SESAT)}}
\end{array}
\]

\begin{contoh}
Perhatikan dan resapi kesalahan fatal berikut:
\begin{itemize}
    \item \textbf{Janji ($p \rightarrow q$)}: "Jika kompor di dapur meledak, maka seluruh gedung apartemen ini ludes terbakar." (Sangat Benar).
    \item \textbf{Fakta Tragedi ($q$)}: "Di siaran TV Breaking News, Apartemen kita disiarkan ludes terbakar jadi abu."
    \item \textbf{Kesimpulan Sesat ($\therefore p$)}: "Astaga! Pasti Kompor pemicunya yang meledak mengawali ini semua."
\end{itemize}
\textbf{Mengapa cacat?} Apartemen terbakar ($q$) dapat disebabkan banyak hal selain kompor (petir, korsleting, dll). Kebenaran $q$ tidak menggaransi kebenaran $p$.
\end{contoh}

\subsection{2. Menyangkal Anteseden (Fallacy of Denying the Antecedent)}

Fallacy ini merupakan versi keliru dari Modus Tollens.

\begin{teorema}[Jebakan Menyangkal Sebab]
Dari implikasi ($p \rightarrow q$) dan fakta bahwa anteseden salah ($\neg p$), seseorang keliru menyimpulkan konsekuen juga salah ($\neg q$). \textbf{Argumen invalid}.
\end{teorema}

\textbf{Anatomi Argumen Sesat:}
\[
\begin{array}{cl}
p \rightarrow q & \text{(Premis Mayor)} \\
\neg p & \text{(Sang Penulis Sebab digugurkan)} \\
\hline
\therefore \neg q & \text{\xcancel{(KESIMPULAN CACAT! SESAT)}}
\end{array}
\]

\begin{contoh}
\begin{itemize}
    \item \textbf{Janji ($p \rightarrow q$)}: "Jika saya hari ini diguyur gajian bonus akhir tahun miliaran, maka tabungan bank BCA saya meluber lubis kaya raya."
    \item \textbf{Fakta Kemiskinan ($\neg p$)}: "Sial, pak bos tidak memberi saya secuil pun gaji bonus."
    \item \textbf{Kesimpulan Sesat ($\therefore \neg q$)}: "Ya sudah fix, berarti saya ditakdirkan selamanya miskin, tabungan BCA saya akan terus mengkerut tidak bisa membesar raya alias gagal gajian!"
\end{itemize}
\textbf{Mengapa cacat?} Kekayaan ($q$) dapat diperoleh melalui banyak cara (warisan, investasi, dll). Ketidakbenaran $p$ tidak mengimplikasikan ketidakbenaran $q$.
\end{contoh}
